\chapter*{Conclusion and Future Work}
\phantomsection
\addcontentsline{toc}{chapter}{Conclusion and Future Work}

\section*{Summary}
\phantomsection
\addcontentsline{toc}{section}{Summary}
This thesis studied knowledge distillation for heterogeneous vision-language
models where teachers vary in usefulness across samples and differ from the
student in tokenizer. Building on cross-tokenizer distillation, optimal transport and
sparse routing
\cite{boizard2024towardscrosstokenizerdistillation,peyre2019computationaloptimaltransport,fedus2022switchtransformersscalingtrillion},
the work introduced \gracefull{} (\gracename{}), a routed multi-teacher
framework that combines three main contributions: Byte-Trie Wasserstein
tokenizer-agnostic distillation, sparse top-\(k\) teacher routing and gradient-agreement
refinement.

Byte-Trie Wasserstein distillation maps decoded teacher and student tokens
into a shared byte-trie structure and compares subtree probability masses
across tokenizer vocabularies. The objective-aware refinement combines sparse
routing, which selects teachers relevant to each sample, with gradient-agreement
reweighting by the cosine alignment between each teacher's KD gradient and the
supervised answer gradient. The empirical evaluation on TextVQA, DocVQA and
ChartQA shows that this combination outperforms same-family, single-teacher
and heterogeneous multi-teacher baselines.

Across TextVQA, DocVQA and ChartQA, \methodname{} scores highest among the
compared methods in Chapter~\ref{chap:experiment-results}. For
SmolVLM-500M, it reaches 69.38 on TextVQA, 73.11 on DocVQA and 80.20 on
ChartQA. For SmolVLM-256M, it reaches 60.73, 60.48 and 72.32 on the same three
benchmarks. On SmolVLM-500M, \methodname{} closes about 55\% of the original
teacher-student gap on TextVQA and ChartQA and about 40\% on DocVQA.
Closing more than half the gap to teachers six times larger suggests that
the supervision design, rather than model capacity, is the primary bottleneck
for compact VLMs on these benchmarks.

\section*{Limitations}
\phantomsection
\addcontentsline{toc}{section}{Limitations}
Three limitations qualify the current results. \textbf{Cost-limited evaluation.} The reported
single-run design limits uncertainty analysis for small gaps between
high-scoring baselines because confidence intervals and estimates from repeated runs are unavailable. As reported in
Section~\ref{sec:efficiency-analysis}, an average run costs approximately
14.41 USD before offline cache generation and storage cost. Although the main
tables use a shared four-teacher pool and consistent hyperparameters, no uncertainty estimates are available for closely matched methods. Small score differences between high-scoring baselines should be treated with caution until repeated-run estimates are available.

\textbf{Scope of validation.} The empirical study focuses on
VQA benchmarks and SmolVLM-family students. The results support VQA
distillation with compact projector-based students, while broader multimodal
task families, other student architectures and larger teacher pools require
further validation. Tasks with long-form answers, multi-image inputs, video reasoning or
open-ended dialogue require different forms of teacher selection and answer
supervision.

\textbf{Reasoning-oriented model support remains outside the current scope.}
GRACE currently distills teacher supervision through output distributions at answer tokens
and gradient agreement with the supervised answer loss. This design fits the evaluated VQA setting, where the target is a short final
answer. It focuses on supervision of final answers and
leaves intermediate deliberation traces, hidden reasoning tokens and extended
test-time reasoning for future work. If a teacher produces only a final answer
distribution, GRACE can still use that distribution as a teacher signal. Future
extensions should separate tokens in the reasoning phase from those in the
answer phase, route teachers by the quality of intermediate reasoning, and
define distillation objectives for multi-step reasoning traces.

\section*{Future Work}
\phantomsection
\addcontentsline{toc}{section}{Future Work}
Four directions remain. The most immediate gap is statistical: repeated runs
with multiple random seeds would clarify which score differences between close
multi-teacher baselines are reliable. On the systems side, more compact logit
storage, adaptive cache loading and lower-cost agreement scoring would reduce
the cost of scaling the method to larger datasets and teacher pools. Beyond
output-level supervision, visual token alignment, cross-attention transfer and
region-level matching could complement the current KD signal for tasks where
the student must localize evidence before answering. Reasoning-oriented models
raise a harder problem: distinguishing reasoning-phase tokens from
final-answer tokens and specifying how teachers are selected when intermediate
traces are available. Router input features, principled conflict resolution and
teacher pools that combine specialized document, chart, general VLM and
reasoning-oriented models are additional open questions
\cite{yuan2020reinforcedmultiteacherselection,zhang2023adaptivemultiteachermetalearning,yu2024selectdistill}.
